\section{Zigzag (Spiral) Level Order Traversal}
\label{sec:trees-zigzag}

\problemstatement{Given the root of a binary tree, return its zigzag level
order traversal: the first level left-to-right, the second right-to-left,
the third left-to-right, and so on, alternating direction each level.}

\begin{patternbox}
\emph{``Zigzag,'' ``spiral,'' ``alternate direction each level''} is plain
\textbf{level-order BFS} with one extra bit of state: a flag that flips
every level and controls whether the current level is emitted forwards or
backwards. The traversal structure does not change at all.
\end{patternbox}

\subsection*{Understanding the problem carefully}

BFS already produces nodes grouped by level, left to right. To zigzag, we
only need to reverse the ordering of \emph{every other} level before
adding it to the answer. The cleanest way is to keep the standard
left-to-right BFS but, when the direction flag says ``right-to-left,''
place each node's value into its position from the back of the current
level's list.

\begin{ideabox}[Standard BFS, Reverse Alternate Levels]
Do an ordinary level-order traversal. Maintain a boolean
\textit{leftToRight}. For each level of size $s$, allocate a list of size
$s$; when \textit{leftToRight} is true, write the $i$-th dequeued node at
index $i$; when false, write it at index $s-1-i$. Flip
\textit{leftToRight} after each level.
\end{ideabox}

\subsection*{Algorithm}

\begin{enumerate}
  \item If root is null, return empty. Push root; \textit{leftToRight} $=$
        true.
  \item While the queue is non-empty:
  \begin{enumerate}
    \item $s=$ queue size; make a level array of length $s$.
    \item For $i=0\ldots s-1$: pop node; index $=$ \textit{leftToRight}
          $?\ i : s-1-i$; store value at that index; push children.
    \item Append the level array; flip \textit{leftToRight}.
  \end{enumerate}
\end{enumerate}

\subsection*{Dry run}

\dryrun{Tree: root $1$; children $2,3$; $2$'s children $4,5$; $3$'s
children $6,7$.}

\begin{itemize}
  \item Level $0$ (L$\to$R): $[1]$.
  \item Level $1$ (R$\to$L): nodes dequeued $2,3$ written at indices
        $1,0 \Rightarrow [3,2]$.
  \item Level $2$ (L$\to$R): $[4,5,6,7]$.
  \item Result: $[[1],[3,2],[4,5,6,7]]$.
\end{itemize}

\subsection*{C++ implementation}

\begin{lstlisting}
#include <bits/stdc++.h>
using namespace std;

struct TreeNode {
    int val;
    TreeNode* left;
    TreeNode* right;
    TreeNode(int v) : val(v), left(nullptr), right(nullptr) {}
};

vector<vector<int>> zigzagLevelOrder(TreeNode* root) {
    vector<vector<int>> ans;
    if (root == nullptr) return ans;

    queue<TreeNode*> q;
    q.push(root);
    bool leftToRight = true;
    while (!q.empty()) {
        int s = q.size();
        vector<int> level(s);
        for (int i = 0; i < s; i++) {
            TreeNode* node = q.front();
            q.pop();
            int idx = leftToRight ? i : (s - 1 - i);
            level[idx] = node->val;
            if (node->left)  q.push(node->left);
            if (node->right) q.push(node->right);
        }
        leftToRight = !leftToRight;
        ans.push_back(level);
    }
    return ans;
}

int main() {
    TreeNode* root = new TreeNode(1);
    root->left = new TreeNode(2);
    root->right = new TreeNode(3);
    root->left->left = new TreeNode(4);
    root->left->right = new TreeNode(5);
    root->right->left = new TreeNode(6);
    root->right->right = new TreeNode(7);

    for (auto& lvl : zigzagLevelOrder(root)) {
        for (int x : lvl) cout << x << " ";
        cout << "\n";
    }
    return 0;
}
\end{lstlisting}

\begin{complexitybox}
\textbf{Time Complexity.} $O(n)$ -- one BFS pass; the index arithmetic is
$O(1)$ per node. \textbf{Space Complexity.} $O(n)$ for the queue and
output.
\end{complexitybox}

\begin{takeawaybox}
Many ``fancy-looking'' traversals are ordinary BFS with a single extra
piece of state controlling how each level is read off. Spot the underlying
level-order engine and the decoration -- here, a direction flag -- becomes
trivial.
\end{takeawaybox}
